\documentclass{article}
\usepackage[utf8]{inputenc}
\usepackage[T1]{fontenc}
\usepackage[spanish, es-tabla]{babel}
\usepackage{booktabs}
\usepackage{multirow}
\usepackage{amsmath}
\usepackage{caption}
\usepackage{makecell} 
\usepackage[legalpaper, margin=1.2cm, top=1.2cm, bottom=1.2cm]{geometry} 

\renewcommand\theadfont{\bfseries} 
\renewcommand\theadalign{cc}       

\begin{document}

\setcounter{table}{7}

% --- TABLA 2: CONFIGURACIÓN DE HIPERPARÁMETROS ---
\begin{table}[ht]
    \centering
    \caption{Configuración de hiperparámetros para experimentos seleccionados (Telco Churn)}
    \label{tab:configuracion_telco_churn}
    \small
    \begin{minipage}{0.85\textwidth} 
        \centering
        \begin{tabular}{cllcccccc}
            \toprule
            \thead{Exp.} & \thead{Umbral IF*} & \thead{Técnica} & \thead{PR\_dist} & \thead{PR\_ries} & \thead{Criterio} & \thead{U\_den} & \thead{U\_ries} & \thead{Tipo\_P} \\
            \midrule
            \textbf{E1} & 1 (1480) & pcsmote & 85 & 30 & proporción & 060 & 040 & Upp060 \\
            \textbf{E2} & 0 (1495) & pcsmote & 85 & 40 & entropía & 060 & 045 & PE70 \\
            \bottomrule
        \end{tabular}
        \vspace{1mm}
        \begin{flushleft}
            \footnotesize * Valor entre paréntesis indica semillas sintéticas candidatas detectadas.
        \end{flushleft}
    \end{minipage}
\end{table}

\vspace{0.4cm}

% --- TABLA 3: RENDIMIENTO EN VALIDACIÓN CRUZADA ---
\begin{table}[ht]
    \centering
    \caption{Rendimiento en Validación Cruzada ($k$-fold CV) para el Caso Base (Telco Churn)}
    \label{tab:cv_telco_churn}
    \small
    \begin{minipage}{0.85\textwidth}
        \centering
        \begin{tabular}{lccccc}
            \toprule
            \thead{Configuración} & \thead{Tamaño Train} & \thead{Tamaño Test} & \thead{F1\_M CV} & \thead{F1\_W CV} & \thead{Acc CV} \\
            \midrule
            Caso Base (Original) & 5634 & 1409 & 0,711 & 0,782 & 0,790 \\
            \bottomrule
        \end{tabular}
        \vspace{1mm}
        \begin{flushleft}
            \footnotesize \textit{Nota.} El conjunto de entrenamiento (80\%) presenta un fuerte desbalance de clases: 4139 instancias pertenecen a la clase mayoritaria (-1) y solo 1495 a la clase minoritaria (1).
        \end{flushleft}
    \end{minipage}
\end{table}

\vspace{0.4cm}

% --- TABLA 4: RENDIMIENTO FINAL ORDENADO POR F1_M TEST ---
\begin{table}[ht]
    \centering
    \caption{Rendimiento final en conjunto de prueba ordenado por F1 Macro (Telco Churn)}
    \label{tab:resultados_telco_churn_ordenados}
    \small
    \setlength{\tabcolsep}{3pt} % Ajuste para la nueva columna
    \begin{minipage}{\textwidth}
        \centering
        \begin{tabular}{cclccccccrc}
            \toprule
            \thead{Ranking} & \thead{Umbral IF} & \thead{Técnica} & \thead{Sem. Val.*} & \thead{Sintet.} & \thead{Train Final\\(Orig + Sintet.)} & \thead{F1\_M\\Test} & \thead{$\Delta$\\(F1\_M)} & \thead{F1\_W\\Test} & \thead{Acc\\Test} \\
            \midrule
            -- & 0\% & \textbf{Caso Base} & -- & -- & 5634 & 0,705 & -- & 0,779 & 0,789 \\
            -- & 1\% & \textbf{Caso Base} & -- & -- & 5577 & 0,702 & -- & 0,776 & 0,786 \\
            \midrule
            1 & \multirow{4}{*}{\shortstack{0\% \\ (candidatas: \\ 1495)}} & smote & -- & 2644 & 8278 & 0,711 & +0,006 & 0,782 & 0,790 \\
            2 & & adasyn & -- & 2594 & 8228 & 0,710 & +0,005 & 0,773 & 0,771 \\
            3 & & \textbf{pcsmote (E2)} & \textbf{897} & \textbf{2644} & \textbf{8278} & \textbf{0,708} & \textbf{+0,003} & \textbf{0,775} & \textbf{0,778} \\
            4 & & borderline & -- & 2644 & 8278 & 0,708 & +0,003 & 0,770 & 0,768 \\
            \midrule
            1 & \multirow{4}{*}{\shortstack{1\% \\ (candidatas: \\ 1480)}} & \textbf{pcsmote (E1)} & \textbf{495} & \textbf{2617} & \textbf{8194} & \textbf{0,715} & \textbf{+0,013} & \textbf{0,783} & \textbf{0,788} \\
            2 & & smote & -- & 2617 & 8194 & 0,713 & +0,011 & 0,776 & 0,776 \\
            3 & & adasyn & -- & 2535 & 8112 & 0,712 & +0,010 & 0,776 & 0,771 \\
            4 & & borderline & -- & 2617 & 8194 & 0,703 & +0,001 & 0,766 & 0,763 \\
            \bottomrule
        \end{tabular}
        \vspace{1mm}
        \begin{flushleft}
            \footnotesize * Semillas candidatas validadas en train. $\Delta$ (F1\_M) representa la mejora absoluta respecto al Caso Base del umbral correspondiente. Nota: Las técnicas están ordenadas por F1\_M Test para mostrar el ranking real de pcsmote.
        \end{flushleft}
    \end{minipage}
\end{table}


\end{document}